\documentclass[a4paper, 12pt]{article}

\usepackage{custom}


%--------------------VARIABILI--------------------
\def\lastversion{v0.1}
\def\title{Verbale interno}
\def\date{24 Marzo 2025}
%------------------------------------------------

\begin{document}

\primapagina


\begin{registromodifiche}
        \lastversion & 26 Marzo 2025 & Pedro Leoni & Marko Peric & Stesura verbale interno \\
        \hline 
\end{registromodifiche}

\tableofcontents

\newpage

\section{Registro presenze}
\begin{itemize}
    \item[] \textbf{Data}: 24 Marzo 2025
    \item[] \textbf{Ora inizio}:  10:00
    \item[] \textbf{Ora fine}: 14:00
    \item[] \textbf{Piattaforma}: Discord	
\end{itemize}

\begin{table}[H]
\centering
{\renewcommand{\arraystretch}{2}
\begin{tabularx}{\textwidth}{| X | X |}
    \hline
        \textbf{\large Componente} & 
        \textbf{\large Presenza} \\ 
    \hline 
    \hline
        Eghosa Matteo Igbinedion Osamwonyi&
        Presente \\
    \hline 
        Guirong Lan&
        Presente \\
    \hline 
        Enrico Bianchi&
        Presente \\
    \hline 
        Francesco Savio&
        Presente \\
    \hline 
        Marko Peric&
        Presente \\
    \hline 
        Pedro Leoni&
        Presente \\
    \hline 

\end{tabularx}}
\end{table}

\newpage

\section{Verbale}

\subsection{Ordine del giorno}
\begin{itemize}
    \item Analisi dei consigli migliorativi dati dai Professori;
    \item Definizione della milestone relativa alla Product Baseline;
    \item Pianificazione delle attività da eseguire fino alla milestone sopracitata;
    \item Riorganizzazione delle risorse per ruolo rimaste in funzione della pianificazione.
\end{itemize}

\subsection{Decisioni prese}
\begin{itemize}
    \item Assegnazione del Responsabile e dei Verificatori in modo fisso indipendentemente dalle attività da portare a termine nello \glossario{sprint}.
    In questo modo il Responsabile può essere più presente e può seguire il lavoro degli altri membri con più attenzione.
    Inoltre si cerca di evitare un problema riscontrato in passato in cui le \glossario{issue} venivano terminate ma non verificate;

    \item Modifica della data della milestone relativa alla Product Baseline al giorno \textbf{30 Maggio 2025}.
    La data inizialmente stabilita (20 Aprile 2025) è stata definita insufficiente, dopo una rivalutazione della pianificazione tenendo in conto lo stato di avanzamento attuale del progetto; 

    \item Modifica delle ore per ruolo preventivate sottraendo ore ai ruoli di Programmatore e Analista in favore dei ruoli di Progettista e Verificatore.
    Il gruppo, dopo uno studio/allenamento individuale sull'attività di progettazione, ha valutato le ore destinate ai ruoli di Progettista e Verificatore come insufficienti.
    Inoltre, ponendo più attenzione alle attività di progettazione e stesura dei test il gruppo ha preventivato di riuscire a eseguire l'attività di codifica in due \glossario{sprint};  

    \item Applicare immediatamente le correzioni indicate dai Professori, così da non accumulare troppi errori;
    
    \item Temporizzare le richieste di modifica indicando data di termine.
    Questa scelta cerca di arginare il problema delle \glossario{issue} irrisolte riscontrato in passato dal team;

    \item Pre assegnamento delle richieste di modifica ai membri del team. 
    Si cerca di assegnare più responsabilità ai singoli membri del team sperando in un miglioramento in termini di qualità dei prodotti e un rispetto maggiore delle tempistiche;

    \item Diminuzione delle dimensioni delle \glossario{issue}, così da evitare di scoprire errori di grandi dimensioni tardi. 
\end{itemize}

\section{To Do/In progress}
\begin{tabularx}{\textwidth}{|c|C|C|C|}
    \hline
        \textbf{Issue} & \textbf{Descrizione} & \textbf{Incaricato} & \textbf{Scadenza}\\
    \hline
        \#265 & Correzione processo di Sviluppo delle Norme di Progetto & Pedro Leoni & 26 Marzo 2025 \\ 
    \hline
        \#266 & Correzione processo di Fornitura delle Norme di Progetto & Marko Peric & 25 Marzo 2025 \\ 
    \hline 
        \#267 & Correzione processo di Documentazione delle Norme di Progetto & Guirong Lan  & 26 Marzo 2025 \\ 
    \hline 
        \#268 & Correzione processo di Gestione della Configurazione delle Norme di Progetto & Francesco Savio & 26 Marzo 2025 \\ 
    \hline
        \#269 & Correzione processo di Verifica delle Norme di Progetto & Guirong Lan & 25 Marzo 2025 \\ 
    \hline
        \#270 & Correzione processo di Accertamento delle Qualità delle Norme di Progetto & Enrico Bianchi & 25 Marzo 2025 \\ 
    \hline
        \#271 & Stesura processo di Risoluzione dei Problemi delle Norme di Progetto & Marko Peric & 25 Marzo 2025 \\ 
    \hline
        \#272 & Stesura processo di Miglioramento delle Norme di Progetto & Eghosa Matteo Igbinedion Osamwonyi & 25 Marzo 2025 \\ 
    \hline
        \#273 & Stesura processo di Gestione delle Norme di Progetto & Enrico Bianchi & 26 Marzo 2025 \\ 
    \hline
        \#274 & Stesura sezioni Strumenti delle Norme di Progetto & Francesco Savio & 26 Marzo 2025 \\ 
    \hline
        \#291 & Stesura processo di Codifica delle Norme di Progetto & Eghosa Matteo Igbinedion Osamwonyi & 26 Marzo 2025 \\ 
    \hline
        \#292 & Stesura attività di Progettazione delle Norme di Progetto & Marko Peric & 26 Marzo 2025 \\ 
    \hline
        \#293 & Stesura attività di Testing delle Norme di Progetto & Enrico Bianchi & 27 Marzo 2025 \\ 
    \hline
        \#294 & Stesura attività di Validazione delle Norme di Progetto & Marko Peric & 27 Marzo 2025 \\ 
    \hline
        \#295 & Stesura metriche legate al prodotto delle Norme di Progetto & Enrico Bianchi & 28 Marzo 2025 \\ 
    \hline
        \#296 & Aggiornamento piano di progetto & Pedro Leoni & 29 Marzo 2025 \\ 
    \hline
        \#297 & Stesura Introduzione documento Specifica Tecnica & Marko Peric & 28 Marzo 2025 \\ 
    \hline
        \#298 & Stesura Verbale Interno & Pedro Leoni & 26 Marzo 2025 \\ 
    \hline
        \#301 & Stesura test di sistema da RFO-1 a RFO-8 & Marko Peric & 29 Marzo 2025 \\ 
    \hline
        \#302 & Stesura test di sistema da RFO-9 a RFO-16  & Francesco Savio & 29 Marzo 2025 \\ 
    \hline
        \#303 & Stesura test di sistema da RFO-17 a RFO-19 e da RFF-1 a RFF-5 & Enrico Bianchi & 29 Marzo 2025 \\ 
    \hline
        \#304 & Stesura test di sistema da RFF-6 a RFF-13 & Eghosa Matteo Igbinedion Osamwonyi & 29 Marzo 2025 \\ 
    \hline
        \#305 & Stesura test di sistema da RFF-14 a RFF-20 & Guirong Lan & 29 Marzo 2025 \\ 
    \hline
 \end{tabularx}

\end{document}